\documentclass[12pt]{article}
\usepackage[T1]{fontenc}
\usepackage[utf8]{inputenc}
\usepackage[brazil]{babel}
\usepackage[margin=1in]{geometry}
\setlength{\parindent}{0pt}
\begin{document}
\section*{Tema 66}
Processo: Processo\\
Situacao: Cancelado - Tema 355/TNU\\
Ramo: DIREITO PREVIDENCIÁRIO\\
Questao: Saber se é possível o reconhecimento da atividade de seminarista como tempo de serviço para fins previdenciários da mesma forma como se dá ao aluno aprendiz.\\
Tese: O tempo de seminarista em congregação religiosa se aproveita para fins previdenciários, desde que atendidos os mesmos pressupostos exigidos do aluno aprendiz de escola pública profissionalizante. Vide Súmula 18/TNU\\
Fonte: NAO\_INFORMADO
\end{document}
